\documentclass[12pt]{article}
\usepackage[utf8]{inputenc}
\usepackage[T1]{fontenc}
\usepackage[margin=1in]{geometry}
\usepackage{setspace}
\usepackage{xcolor}
\usepackage[colorlinks=true, urlcolor=blue!60!black, linkcolor=blue!60!black]{hyperref}
\setlength{\parindent}{0em}
\setlength{\parskip}{0.6em}
\renewcommand{\baselinestretch}{1.0}
\begin{document}
\begin{center}
{\LARGE \textit{Loving v.\ Virginia}, 388 U.S.\ 1 (1967)\par}
\vspace{0.5em}
{\large Case Brief --- Two-Page Edition\par}
\end{center}

\section{Memory Jogger}
Virginia's ban on interracial marriage violated both equal protection and due process (\href{https://supreme.justia.com/cases/federal/us/388/1/}{Justia}).

\section{Detailed Case Facts}
Mildred Jeter, a Black woman, and Richard Loving, a white man, married in Washington, D.C., in 1958 and returned to Caroline County, Virginia, where police arrested them in a night raid. Virginia made interracial marriage a felony and barred couples from marrying elsewhere to evade the ban.

\section{Procedural History}
The Lovings pleaded guilty; their one-year sentences were suspended on condition that they leave Virginia for twenty-five years. The Supreme Court of Appeals of Virginia upheld the statutes. The U.S.\ Supreme Court reversed.

\section{Judicial Votes}
Unanimous. Warren wrote for eight Justices; Stewart concurred in the judgment.

\section{Holding}
Racial classifications in marriage laws, even if ``equally applied,'' get ``the most rigid scrutiny'' and fail absent a purpose independent of invidious discrimination. Independently, the freedom to marry is a fundamental liberty that race cannot be used to deny.

\section{Analysis of Opinions}
\textbf{Majority:} equal application does not immunize a racial classification (repudiating \textit{Pace v.\ Alabama}); the framers' debates were ``at best, inconclusive''; and banning only marriages involving white persons showed the laws were ``designed to maintain White Supremacy.'' Marriage is one of the ``basic civil rights of man.'' \textbf{Stewart:} a law making ``the criminality of an act depend upon the race of the actor'' can never be valid.

\section{Examples: Future Applications}
\textbf{Same-Side:} a ban on marriage across national origins fails strict scrutiny. \textbf{Opposite-Side:} age, licensing, and consanguinity rules are valid police-power regulation. \textbf{Fence-Sitter:} a bar on marriage for parents behind on child support, struck down in \textit{Zablocki v.\ Redhail}.

\section{Critique}
\textit{Loving} is a litmus test for originalism, since many framers expected anti-miscegenation laws to survive; originalists answer with the Amendment's equal-civil-rights text. Others find the due process part under-theorized.

\section{Key Quotations}
\begin{quote}
``The freedom to marry has long been recognized as one of the vital personal rights essential to the orderly pursuit of happiness by free men.''
\end{quote}
\noindent --- Chief Justice Warren, for the Court, 388 U.S.\ at 12.

\end{document}
